\section{Methodology}
\label{sec:methodology}

\subsection{Study Design and Reporting Scope}

This study reports a descriptive offline benchmark of Spanish obstetric question answering. The primary comparison estimates the performance difference between each base model and its grounded QLoRA adapter when both receive supplied source evidence. The comparison uses the same held-out test variant and documented prompt construction, decoding, output linkage, and evaluation procedures for all four conditions. It therefore measures evidence-conditioned answer generation; it does not measure a closed-book parametric adaptation effect. The benchmark is not clinical validation, expert validation, patient-outcome assessment, or evidence for deployment or clinical decision support. Reporting follows transparency principles associated with TRIPOD-LLM without claiming TRIPOD-LLM compliance or clinical validation \cite{Gallifant2025TRIPODLLM}.

Four evaluation conditions were retained, and each condition contains 328 test outputs. Table~\ref{tab:evaluation-conditions} identifies the base model, adapter status, input variant, and retained metric-evaluation errors for each condition.

\begin{table}[htbp]
\centering
\small
\setlength{\tabcolsep}{3pt}
{\scriptsize
\begin{tabular}{@{}p{0.16\linewidth}p{0.18\linewidth}p{0.23\linewidth}p{0.12\linewidth}r r@{}}
\hline
\shortstack{Evaluation\\condition} & \shortstack{Base\\model} & Adapter & \shortstack{Input\\variant} & \shortstack{Test\\outputs} & \shortstack{Evaluation\\errors} \\
\hline
Gemma 4 base & \scriptsize\texttt{google/}\allowbreak\texttt{gemma-4-E2B-it} & No adapter & \texttt{sft\_grounded} & 328 & 1 \\
Gemma 4 QLoRA & \scriptsize\texttt{google/}\allowbreak\texttt{gemma-4-E2B-it} & \scriptsize\texttt{iue-edu/}\allowbreak\texttt{MaternaCare-ES-}\allowbreak\texttt{gemma4-qlora} & \texttt{sft\_grounded} & 328 & 1 \\
MedGemma 1.5 base & \scriptsize\texttt{google/}\allowbreak\texttt{medgemma-1.5-}\allowbreak\texttt{4b-it} & No adapter & \texttt{sft\_grounded} & 328 & 5 \\
MedGemma 1.5 QLoRA & \scriptsize\texttt{google/}\allowbreak\texttt{medgemma-1.5-}\allowbreak\texttt{4b-it} & \scriptsize\texttt{iue-edu/}\allowbreak\texttt{MaternaCare-ES-}\allowbreak\texttt{medgemma-qlora} & \texttt{sft\_grounded} & 328 & 6 \\
\hline
\end{tabular}
}
\caption{Retained evaluation conditions and output accounting. Evaluation errors are metric-evaluation records reported by the retained summaries; they do not indicate missing model predictions.}
\label{tab:evaluation-conditions}
\end{table}

\subsection{Dataset, Provenance, and Split Integrity}

MaternaQA-es comprises 5,727 question--answer pairs. The training partition contains 5,093 pairs from 52 source PDFs, the validation partition contains 306 pairs from 2 PDFs, and the held-out test partition contains 328 pairs from 3 PDFs. Splitting was performed at the document level, and no test PDF overlaps with the training or validation partitions. This design limits document-level leakage between development and test material, although topical or linguistic similarity across documents remains possible. Dataset documentation identifies provenance, intended use, and limitations in the spirit of data statements \cite{Bender2018DataStatements}.

The primary training and evaluation input is the held-out \texttt{sft\_grounded} variant. It contains source context together with the question; in the conversational record, the user prompt contains \texttt{Contexto fuente} followed by the question. Thus, for each test item, the evidence passage is supplied as part of the dataset record before inference and is not dynamically retrieved by a RAG component during these four evaluations. The \texttt{sft\_closed\_book} variant contains the question only and exists as a separate dataset variant, but it is not part of the primary comparison. Here, \emph{base} means the base weights without an adapter; it does not mean closed-book inference. Accordingly, the primary results estimate base-versus-grounded-QLoRA performance under supplied evidence rather than a closed-book parametric adaptation effect.

\subsection{Model Families and QLoRA Adaptation}

The evaluated base model families are \texttt{google/}\allowbreak\texttt{gemma-4-E2B-it} and \texttt{google/}\allowbreak\texttt{medgemma-1.5-4b-it}. The corresponding retained adapters are \texttt{iue-edu/}\allowbreak\texttt{MaternaCare-ES-}\allowbreak\texttt{gemma4-qlora} and \texttt{iue-edu/}\allowbreak\texttt{MaternaCare-ES-}\allowbreak\texttt{medgemma-qlora}. QLoRA adapts a quantized language model through low-rank updates \cite{Dettmers2023QLoRA}. The training configuration uses a 4-bit NF4 backbone with double quantization and LoRA over all linear layers. The retained trainer states confirm two epochs and 1,274 training steps for both grounded adapters.

Context-conditioned supervised fine-tuning is related in motivation to retrieval-augmented adaptation, including RAFT \cite{Zhang2024RAFT}, but the present procedure is not a RAFT implementation. The primary comparison uses supplied source context in the supervised input and does not claim to reproduce RAFT's distractor-document or chain-of-thought components.

\subsection{Training and Inference Protocol}

Training was performed with \texttt{scripts/train\_qlora\_trl.py}. The script supports loading JSON datasets from local files or Hugging Face, and it loads separate training, validation, and held-out test splits. Each conversational example is divided into a prompt and completion at the final assistant message. The retained grounded-adapter configuration used 4-bit NF4 quantization with double quantization; LoRA over all linear layers with rank $r=16$, alpha $=16$, dropout $=0.05$, and bias set to \texttt{none}; completion-only loss; maximum sequence length 1024; two epochs; per-device batch size 1; gradient accumulation of 8; learning rate $2\times10^{-4}$; warmup ratio 0.03; weight decay 0; maximum gradient norm 0.3; 8-bit AdamW; a cosine scheduler; and seed 3407. The retained trainer states confirm two epochs and 1,274 steps for both grounded adapters.

Both base and QLoRA inference scripts construct the prompt from the same message structure. They remove the reference assistant message, render the remaining user and system messages with the tokenizer chat template using \texttt{add\_generation\_prompt=True}, tokenize the rendered prompt, generate a completion, and decode only the newly generated tokens. Consequently, the base and QLoRA conditions receive the same question and the same source-context field for each paired \texttt{qa\_id}; the experimental difference is the presence or absence of the trained adapter, not a difference in retrieved evidence. The retained output records preserve \texttt{qa\_id}, question, source context, generated answer, reference answer, rendered prompt text, model role and name, adapter directory, and dataset variant. These fields maintain the linkage between each output and its test item.

The documented inference defaults are \texttt{max\_new\_tokens}=512 and \texttt{do\_sample}=False, which specifies greedy deterministic decoding, together with the tokenizer's \texttt{pad\_token\_id} and \texttt{eos\_token\_id}. Temperature 0.7 and top-$p$ 0.9 are supplied only when sampling is explicitly enabled. The QLoRA script also supports optional repetition-penalty and no-repeat-ngram controls. Those controls are not claimed to have been active in the retained original results.

\subsection{Outcome Measures and Analysis Plan}

The retained evaluation used the \texttt{ragas==0.4.3} stack. The judge was \texttt{gpt-5.4-mini}, the embedding model was \texttt{text-embedding-3-small}, and the maximum completion length was 2048 tokens. Faithfulness evaluates the generated answer against the supplied source context. Answer Relevancy evaluates the relevance of the generated answer to the question. Answer Correctness evaluates the generated answer against the reference answer. Semantic Similarity evaluates embedding similarity between the generated and reference answers \cite{Es2023Ragas}.

The retained summaries contain 1, 1, 5, and 6 metric-evaluation errors for Gemma 4 base, Gemma 4 QLoRA, MedGemma 1.5 base, and MedGemma 1.5 QLoRA, respectively. These are records of evaluation failures or unavailable metric values, not absent predictions. The retained analysis is descriptive; inferential statistics are outside its scope. If a later rerun adds inferential analysis, it should pre-specify the estimand, metric family, resampling procedure, random seed, interval method, multiplicity handling, and treatment of metric-evaluation errors. Any such analysis would quantify uncertainty over the empirical 328-item benchmark distribution, not over new source documents.

\subsection{Auxiliary Retrieval-Augmented Evaluation}

Auxiliary RAG is separate from grounded supervised fine-tuning. The current repository implements \texttt{no\_rag}, Hybrid retrieval combining BM25 and dense retrieval with reciprocal-rank fusion (RRF) \cite{Cormack2009RRF}, and independent HyDE, which uses a dedicated hypothetical-document generator followed by dense retrieval \cite{Gao2023HyDE}. These components belong to an auxiliary retrieval-augmented protocol in the broader RAG framework \cite{Lewis2020RAG}; they are not additional conditions in the four retained fine-tuning summaries. The RAGAS evaluation framework is cited separately for the metric stack \cite{Es2023Ragas}.

Auxiliary RAG results should be reported only when the corresponding corpus, index, retrieval, generator, configuration, and evaluation artifacts are available. The four retained fine-tuning summaries must not be presented as Hybrid, HyDE, or other RAG results. Grounded supervised fine-tuning supplies source context as part of the model input, whereas auxiliary RAG retrieves evidence at inference time; the two protocols therefore answer different methodological questions.

\subsection{Threats to Validity and Scope}

Several limitations constrain interpretation. The QA pairs are synthetic and may paraphrase, simplify, omit qualifications, or overgeneralize the source material. The source scope is limited to the documents represented in MaternaQA-es, and only three PDFs are held out for testing. Automated judge and embedding-based metrics provide useful signals but cannot capture every factual, clinical, or safety error; the retained metric-evaluation errors add a further limitation to metric completeness. No clinical or expert validation, patient-outcome assessment, deployment evaluation, or decision-support assessment was performed.

The primary results also do not include a closed-book ablation or a counterfactual context ablation. They therefore do not isolate a parametric adaptation effect without supplied evidence. The conclusions are restricted to the descriptive offline comparison of the specified base models and grounded QLoRA adapters on the retained 328-item test set. Future clinical or deployment-oriented evaluation would require a separately specified protocol, governance review, and evidence beyond this benchmark.
